\chapter{Release 4 --- Consolidation, Delivery and Deployment}

\section*{Introduction}
The functional scope of the platform was completed at the end of Release~3. This final release makes
it deliverable: a coherent and accessible user interface, a verified authorization policy, a
security review, and a reproducible deployment. It gathers two sprints: Sprint~8 consolidates the interface and
verifies the platform, Sprint~9 makes it reproducibly deployable and continuously verified.
Neither delivers new business functionality; both deliver deliverability.

% =====================================================================
\section{Sprint 8 --- Interface Consolidation, Testing and Deployment}

\subsection{Sprint Objective}
Bring the platform to a deliverable state: unify the user interface, verify that the security model
behaves as designed, review the application against the common web vulnerabilities, and package it
for deployment.

\subsection{Functional Scope}
This sprint covers user story US-26, a consistent and responsive interface, and the
cross-cutting quality work that no single functional story carries.

\subsection{Features Implemented}
\begin{itemize}
  \item Frontend application structure and session handling.
  \item Authorization model applied to the interface.
  \item Design system unifying the screens of the eight modules.
  \item Automated test suite covering the business rules and the access control.
  \item Security review against the common web vulnerabilities.
  \item Deployment packaging.
\end{itemize}

\subsection{Technical Realisation}

\subsubsection{Frontend structure and session handling}
The client is a single-page application organised by feature, mirroring the backend modules. The
access token obtained at sign-in is kept in memory and attached to every request by an HTTP
interceptor; the refresh token remains in its \texttt{HttpOnly} cookie, unreachable from
JavaScript. When a request returns \texttt{401}, the interceptor attempts a single refresh before
returning the user to the sign-in screen.

\subsubsection{Authorization applied to the interface}
The interface consumes the permission list returned by \texttt{GET /api/me/context} and derives
from it both the navigation menu and the visibility of individual actions. A button that the user
is not entitled to use is not merely disabled: it is absent.

This must be understood correctly, and the report states it explicitly: hiding a control is an
\textbf{ergonomic} measure, not a security measure. The authorization decision is taken by the
server on every request, as described in Sprint~2. The client-side rule exists so that the user is
never offered an action that would be refused; it is never the thing that refuses it.

\subsubsection{Design system and interface consolidation}
The eight modules were built successively, each with its own screens. This sprint unified them
behind a single design system: a token-based visual language (colours, typography, spacing,
elevation) declared once and consumed by every screen, together with a set of shared components
(pagination, selection, confirmation dialog, notifications) reused across the modules rather than
reimplemented in each.

The consolidation also addressed accessibility and consistency: sortable table headers announcing
their state, labelled icon-only controls, explicit empty states in place of blank tables, and a
responsive layout. The interface being the only part of the platform its users ever see, this
convergence matters as much to adoption as the correctness of the engine behind it.

\subsubsection{Test strategy}
Testing concentrates on what a regression would make expensive: the business rules and the
authorization model. The automated backend suite covers the lifecycle transitions, the billing
ceiling, workload uniqueness and immutability, the multi-currency conversions, and the indicator
formulas compared against the reference workbook.

The access control receives dedicated attention: for each protected endpoint, the suite verifies
that a caller holding the permission succeeds and that a caller without it is refused. This is what
makes the ``configurable authorization'' claim verifiable rather than merely stated.

\subsubsection{Security review}
The application was reviewed against the common categories of web vulnerability:

\begin{longtable}{|m{4.6cm}|m{10cm}|}
\hline
\textbf{Category} & \textbf{Measure taken} \\
\hline
\endhead
\endfoot
\endlastfoot
\hline
Broken access control & Every endpoint annotated with its required permission; scope checked
independently of permission; no role name tested in code \\
\hline
Cryptographic failures & Passwords hashed with BCrypt; tokens signed; refresh token confined to an
\texttt{HttpOnly}, \texttt{Secure}, \texttt{SameSite} cookie \\
\hline
Injection & Parameterised queries through the ORM; no string-built SQL \\
\hline
Identification and session failures & Short-lived access token; instant revocation by token
version; forced password change on first login \\
\hline
Sensitive data exposure & Financial data filtered by permission; internal quote behind a dedicated
permission; computed amounts never persisted \\
\hline
Security misconfiguration & Schema managed by versioned migrations with the ORM in validation
mode; secrets externalised as environment variables \\
\hline
\caption{Security review by vulnerability category}
\label{tab:owasp}
\end{longtable}

\subsubsection{Packaging}
At the close of this sprint the platform is packaged as an executable archive for the backend and a
static build for the frontend, with the database schema brought up to date by the migration tool at
startup and all configuration supplied by environment variables.

Startup is idempotent: deploying the same artefact against an up-to-date database changes nothing,
and against an older one applies exactly the missing migrations.

This packaging is sufficient to run the application, but not to \emph{deploy} it reproducibly:
it still assumes a host on which the right Java version, the right Node version and a correctly
provisioned PostgreSQL instance are already present. Sprint~9 removes that assumption.

\subsection{Tests and Validation}

\begin{longtable}{|m{4.2cm}|m{5.4cm}|m{4.8cm}|}
\hline
\textbf{Test case} & \textbf{Steps} & \textbf{Expected result} \\
\hline
\endhead
\endfoot
\endlastfoot
\hline
Business rule suite & Run the automated backend tests & All pass \\
\hline
Permission granted & Call each protected endpoint with the permission & Operation performed \\
\hline
Permission missing & Call each protected endpoint without it & HTTP 403 \\
\hline
Financial confidentiality & Browse the interface as a developer & No financial figure exposed \\
\hline
Menu consistency & Sign in with each of the four roles & Each menu matches the permissions held \\
\hline
Deployment & Start the artefact against an empty and an existing database & Schema created, then
migrated without loss \\
\hline
\caption{Test cases --- Sprint 8}
\label{tab:tests-s8}
\end{longtable}

\subsection{Sprint Outcome}
The platform is functionally complete, verified on its business rules and its authorization model,
reviewed for the common vulnerabilities, and packaged as a runnable artefact.

\subsection{Transition to the Next Sprint}
One gap remains between ``it runs on my machine'' and ``it can be deployed'': the artefact still
depends on a correctly prepared host. Sprint~9 closes that gap with containerization and a
continuous integration pipeline.

% =====================================================================
\section{Sprint 9 --- Containerization and Continuous Delivery}

\subsection{Sprint Objective}
Make the platform reproducibly deployable and continuously verified: package each component as a
container image, orchestrate the whole stack with a single command, and automate build, test and
integration checks on every change.

\subsection{Functional Scope}
This sprint delivers no new business functionality. Its scope is the delivery chain itself, and it
addresses an item already open in the project's engineering backlog: the absence of deployment
documentation and of any reproducible deployment procedure.

\subsection{Features Implemented}
\begin{itemize}
  \item Container image for the backend, built in two stages and running as a non-root user.
  \item Container image for the frontend, serving the production build and reverse-proxying the API.
  \item Orchestration of the three services with health checks and startup ordering.
  \item Continuous integration pipeline: build, tests, image construction, end-to-end smoke test.
  \item Deployment documentation and an architecture decision record superseding the previous one.
\end{itemize}

\subsection{Technical Realisation}

\subsubsection{A reconsidered architecture decision}
The project had initially recorded a decision to run natively, without containers, justified by
iteration speed on an already-provisioned machine. That decision served development well but left
two gaps: no reproducible deployment, and a production profile that had never actually been
executed.

The decision was therefore superseded rather than silently contradicted: a new record documents the
move to containers, states what it changes, and confirms what it does \emph{not} change. In
particular, three containers are \textbf{not} three microservices: the modular monolith decision
stands. What is containerized is one application, its database and its static server.

\subsubsection{Image construction}
Both images are built in two stages, so that the build toolchain never reaches the runtime image:
the backend image contains a Java runtime and the application archive, with neither Maven nor
sources; the frontend image contains a web server and the compiled assets, with no Node runtime.
The backend container runs as a dedicated unprivileged user rather than as root.

Tests are deliberately \textbf{not} executed during image construction. Building an image is a
packaging step; the quality gate belongs to the pipeline. Running the suite in both places would
double build time without producing any additional signal.

\subsubsection{Single origin through the reverse proxy}
The web server serves the application and forwards \texttt{/api} to the backend container. The
browser therefore contacts a single origin, which has a concrete consequence: the cross-origin
configuration required by the development setup becomes unnecessary, and the refresh cookie,
restricted to its own path and marked \texttt{SameSite=Strict}, works without any exception.
The containerized deployment is, on this point, simpler and safer than the development loop.

\subsubsection{Startup ordering and health probes}
Composing three services raises an ordering problem: a container being \emph{started} is not the
same as a container being \emph{ready}. The stack therefore relies on health probes rather than on
declaration order. The database is probed over TCP, because the temporary server used during
initialisation accepts only local socket connections and would otherwise be reported ready too
early. The backend exposes a health endpoint, deliberately left unauthenticated but returning no
detail, and is given a generous startup window because its first run applies the full set of
schema migrations, including the demonstration dataset.

\subsubsection{Defects revealed by the exercise}
Preparing the pipeline exposed three pre-existing defects that native execution had concealed. The
most consequential concerned the frontend build configuration: the production configuration never
substituted the production environment file, so the production bundle embedded the development API
address. On the development machine the application appeared to work, because the bundle was calling the
local backend, while it would have been broken on any other host. The second was a build wrapper
documented in the project files but absent from the repository. The third was an ignore rule that
silently excluded the configuration template from version control, precisely the file the
deployment procedure requires.

That these defects surfaced only when the delivery chain was automated is itself the argument for
automating it.

\subsubsection{Pipeline}
The pipeline separates what must always hold from what is only meaningful on integration branches.
Compilation, the backend test suite, the production frontend build and the construction of both
images run on every proposed change: a component that only breaks after merging is not a gate. The
end-to-end test, which starts the full stack and performs a real authentication through the
reverse proxy, runs on integration branches, where its cost is justified. Image publication is
reserved for the main branch.

The frontend job also carries an explicit guard against the defect described above: it fails if the
development address reappears in the production bundle.

\subsection{Tests and Validation}

\begin{longtable}{|m{4.2cm}|m{5.4cm}|m{4.8cm}|}
\hline
\textbf{Test case} & \textbf{Steps} & \textbf{Expected result} \\
\hline
\endhead
\endfoot
\endlastfoot
\hline
Environment substitution & Inspect the production bundle & The development address is absent \\
\hline
Unprivileged execution & Inspect the identity inside the backend container & Non-root user \\
\hline
Full stack startup & Start the three services & All three reported healthy \\
\hline
Schema migration & Count the applied migrations on a pristine volume & All migrations applied \\
\hline
End-to-end authentication & Authenticate through the reverse proxy & Token returned and refresh
cookie set \\
\hline
Client-side routing & Reload a deep application route & The application is served, not a 404 \\
\hline
Persistence & Stop and restart the stack & Data preserved; a volume reset replays the migrations \\
\hline
Development loop & Run the native backend and frontend while the stack is up & Both coexist, no
port collision \\
\hline
\caption{Test cases --- Sprint 9}
\label{tab:tests-s9}
\end{longtable}

\subsection{Sprint Outcome}
The platform is deployable by a single command on any machine running a container engine, and every
change is now verified automatically. The deployment procedure, the environment variables, the
backup method and the known pitfalls are documented.

\subsection{Final State of the Product}
At the end of the nine sprints, PMS covers the whole scope identified during the analysis of the
Excel workbook: dynamic authorization, project and team management, workload planning and
recording, billing, missions, governance and the earned-value indicator engine, all behind a
unified interface. Table~\ref{tab:final} summarises the path from the reference workbook to the
delivered platform.

\begin{longtable}{|m{2.0cm}|m{5.6cm}|m{6.8cm}|}
\hline
\textbf{Sprint} & \textbf{Increment delivered} & \textbf{Excel logic replaced} \\
\hline
\endhead
\endfoot
\endlastfoot
\hline
Sprint 1 & Technical foundation, authentication & --- (no spreadsheet equivalent) \\
\hline
Sprint 2 & Dynamic authorization, users, resources, cost rates & TCC sheet \\
\hline
Sprint 3 & Projects and lifecycle & Identification sheet \\
\hline
Sprint 4 & Teams, workload plan and actuals & Forecast cost and actual cost sheets \\
\hline
Sprint 5 & Billing, payments, amendments, missions & Billing progress and missions sheets \\
\hline
Sprint 6 & Governance registers & Risks, deliverables, stakeholders, changes sheets \\
\hline
Sprint 7 & Internal quote and indicator engine & Statistics and dashboard sheets \\
\hline
Sprint 8 & Interface consolidation, tests, security review & --- (quality) \\
\hline
Sprint 9 & Containerization and continuous delivery & --- (delivery chain) \\
\hline
\caption{From the reference workbook to the delivered platform}
\label{tab:final}
\end{longtable}

\section*{Conclusion}
Release~4 closed the construction of the platform by addressing what functionality alone does not
deliver: a coherent interface, a verified security model and a reproducible deployment. The
incremental path followed since Release~1 (foundation, operations, finance, consolidation)
produced a platform whose every increment was demonstrable in isolation, and whose final state
covers the entire business logic previously carried by the Excel workbook.
